\begin{abstract}
Practitioners routinely craft increasingly verbose prompts under the
assumption that more context yields better responses.
We test this assumption with a controlled study spanning 21,858 evaluated
responses across eight large language models (1B to frontier scale), six
tasks, seven strictly additive prompt-elaboration levels, and eight
supplementary experiments.
Each level appends exactly one elaboration mechanism, so any quality change
between adjacent levels is attributable to that single mechanism.
To avoid resting our conclusions on a single evaluator, we score the four
tasks with objective ground truth (classification, product extraction, QA,
math) using \emph{deterministic} metrics as the primary measure, and use an
LLM judge only where no clean heuristic exists, cross-checked by an independent
second judge (gemini-2.0-flash, 3,915~re-evaluations).
Our central finding is that quality gain for structured tasks
\emph{concentrates at a single schema-resolving layer}: the layer that
communicates the output format or label space.
On the deterministic metric this layer accounts for 87\% of the total gain
for classification (the task label) and 92\% for product extraction (the JSON
format specification, 7/7 models significant, $R^2 > 0.99$). On a
\emph{held-out} task chosen after the main study, it accounts for 95\% for
named-entity recognition (4/5 significant).
The held-out task confirms the mechanism predictively rather than merely
describing it.
Two controls establish that this is driven by information content, not token
count: a randomised layer-ordering ablation (2,986~runs) shows the saturation
\emph{shape} survives reordering while the token threshold moves with the
schema layer's position, and an irrelevant-padding control (1,470~runs) shows
matched-token filler produces flat or declining quality (1/26 conditions
significant).
A hard-example stress test on public benchmarks (multi-step GSM8K, multi-hop
HotpotQA, and multi-constraint instructions) shows that QA insensitivity is
genuine rather than a ceiling artefact, while instruction-following sensitivity
strengthens $6.6\times$ from easy to hard examples. This refines the picture
into three mechanism-anchored classes:
\emph{schema-gated} (classification, extraction, NER), \emph{difficulty-gated}
(instruction following), and \emph{knowledge-gated} (QA, math, summarisation,
at ceiling from the bare input).
Repeated temperature-zero runs confirm these effects exceed run-to-run
sampling noise by 20--139$\times$.
The practical guidance is concrete: for structured tasks, invest the prompt
budget in one well-formed schema specification and stop; for QA and math, the
bare question suffices.
\end{abstract}
\begin{keywords}
prompt engineering; LLM evaluation; prompt saturation; information sufficiency;
schema specification; prompt sensitivity
\end{keywords}
